\section*{AI methodology}
\addcontentsline{toc}{section}{AI methodology}
We developed the approach to the main theorem in summer 2025, without AI assistance.
By summer 2026, we had a draft of the present paper, and the companion paper \cite{md-abmw} was nearly complete.
The draft already had essentially the current proof structure.
It contained a complete proof of the result of \Cref{part:types}, and definitions and lemma statements for the other parts, though some of their proofs were still sketches.
Believing that these sketches could be completed, we used GPT-6 to help check the argument.

In \Cref{part:definable-construction-sequences}, we had proposed a proof of a stronger version of \Cref{cl:from-merge-to-construction}, in which the number of unary predicates is bounded by a constant independent of the length of the input merge sequence.
This would have let us apply the results of \cite{md-abmw} without controlling the sequence length.
The model identified a gap in this proof.
We expected the gap to be easy to repair with the methods already in the paper, but the model did not find a repair.

We therefore asked a narrower question: are there construction sequences of width $n^{o(1)}$ that are locally definable using a constant number of unary predicates?
The model observed that a slightly weaker width bound of $n^\eps$, for any fixed $\eps>0$, can be achieved by tuning the parameter $q$ in Theorem 4.3 in \cite{md-abmw}, so that it produces chains of length $O(r/\eps)$.
Since the number of unary predicates depends only on the length of the input sequence, this gives a bound independent of the size of the graph.
We immediately recognized that this was enough to close the gap within our original proof strategy.

The text of \Cref{part:types,part:definable-construction-sequences} was written by the authors.
We used LLMs to help write some proofs in \Cref{part:mc} that are mechanical adaptations of proofs from \cite{locality-note}.
GPT-6 also proposed a shorter proof of \Cref{lem:ball-types} in \Cref{part:types}, which we verified and adopted.
Finally, we used LLMs to streamline the presentation.
The models used were GPT-6, Claude Opus~5.5 and Claude Fable~5.1.

The authors verified the correctness and originality of all content, including references.
